\chapter{结果与分析}

\section{变量质量检查}

示例数据的质量检查显示，图像类变量对光照和遮挡更敏感，人工调查指标的波动则更多来自测量间隔和记录方式。对这些差异进行显式说明，有助于读者判断后续融合结果的可信度。

\section{融合结果}

融合得分能够同时反映冠层结构和环境响应。图 \ref{fig:jlau-demo-flow} 展示了从观测数据到融合评价的基本流程。

\begin{figure}[htbp]
  \centering
  \fbox{\parbox[c][4.2cm][c]{0.78\textwidth}{\centering 多源观测\\特征对齐\\融合评价}}
  \caption{表型数据融合流程示意}
  \label{fig:jlau-demo-flow}
\end{figure}

\section{讨论}

融合方法的价值在于减少单一指标带来的解释偏差，但其有效性依赖于清晰的数据来源记录和稳定的质量控制规则。对于真实论文，应补充试验重复、模型参数、敏感性分析和外部验证结果。

\section{小结}

本章展示了结果与分析部分的基础写法。用户可在保持标题层级和图表编号规则不变的前提下，替换为真实实验结果。
